%%%%%%%%%%%%%%%%%%%%%%%%%
\spellentry{Explosive Runes}
%%%%%%%%%%%%%%%%%%%%%%%%%

\tagline{Abjuration [Force]}
\begin{description*}
\item[Level:] Sor/Wiz 3
\item[Casting Time:] 1 standard action
\item[Range:] Touch
\item[Duration:] Permanent until discharged (D)
\item[Saving Throw:] See text
\item[Spell Resistance:] Yes
\item[Components:] V, S
\end{description*}

You trace these mystic runes upon a book, map, scroll, or similar object bearing
written information. The \textit{runes} detonate when read, dealing 6d6 points
of force damage. Anyone next to the \textit{runes} (close enough to read them)
takes the full damage with no saving throw; any other creature within 10 feet of
the \textit{runes} is entitled to a Reflex save for half damage. The object on
which the \textit{runes} were written also takes full damage (no saving throw).
You and any characters you specifically instruct can read the protected writing
without triggering the \textit{runes}. Likewise, you can remove the \textit{runes
}whenever desired. Another creature can remove them with a successful \textit{dispel
magic} or \textit{erase} spell, but attempting to dispel or erase the \textit{runes
}and failing to do so triggers the explosion.
Note: Magic traps such as \textit{explosive runes} are hard to detect
and disable. A rogue (only) can use the Search skill to find the \textit{runes
}and Disable Device to thwart them. The DC in each case is 25 + spell level, or
28 for \textit{explosive runes}.